\section{The Boy Catechist and His Martyr Pupils}

With the persecution still burning, catechesis in Alexandria had collapsed:
``there was no one at Alexandria to give instruction in the faith, as all
were driven away by the threat of persecution.'' Pagans began coming to the
young grammar teacher ``to hear the word of God''---the first, Plutarch, died
a martyr; the second, his brother Heraclas, would one day be bishop of
Alexandria (\work{Church History} 6.3.1--2). ``He was in his eighteenth year
when he took charge of the catechetical school'' (\work{Church History}
6.3.3); Demetrius subsequently confirmed what danger had improvised, so that
``the catechetical school had been entrusted to him alone by Demetrius, who
presided over the church'' (\work{Church History} 6.3.8).

The teaching was carried on at the edge of the tribunal. Origen attended the
arrested ``not only\ldots\ while in bonds, and until their final
condemnation,'' but to the place of execution, where he ``with great boldness
saluted the martyrs with a kiss,'' repeatedly escaping the mob himself;
soldiers were posted around the house where he taught, and ``he removed from
house to house'' as pupils kept coming (\work{Church History} 6.3.3--6,
6.4--5). Several of his catechumens---Plutarch, Serenus, Heraclides, Hero,
a second Serenus, the catechumen Herais, and the celebrated Potamiaena---died
in the persecution; Eusebius counts them as the school's proper crown.

Having decided that ``the teaching of grammatical science'' was
``inconsistent with training in divine subjects,'' Origen gave up his
school and his library together: ``he disposed of whatever valuable books of
ancient literature he possessed, being satisfied with receiving from the
purchaser four oboli a day'' (\work{Church History} 6.3.8--9). On that
laborer's pension he built the regime that made him famous: fasting, little
sleep, and that on the floor; no second coat, no shoes ``for a number of
years''; years of abstention from wine; poverty pressed ``to the very
extreme,'' to the point that friends were grieved and his health endangered
(\work{Church History} 6.3.9--12). Eusebius's summary became proverbial:
``his manner of life was as his doctrine, and his doctrine as his life''
(\work{Church History} 6.3.7). When the crowds grew beyond one man, he
divided the school, keeping the advanced and entrusting beginners to Heraclas
(\work{Church History} 6.15).

Two further formations belong to these Alexandrian years. He heard the
lectures of Ammonius---so the pagan Porphyry, who despised the result,
attests: Origen ``having been a hearer of Ammonius, who had attained the
greatest proficiency in philosophy of any in our day,'' turned Greek learning
to the service of ``foreign fables,'' ``in his life conducting himself as a
Christian and contrary to the laws, but in his opinions of material things
and of the Deity being like a Greek'' (Porphyry, \work{Against the
Christians}, quoted in \work{Church History} 6.19.5--7). Whether this
Ammonius is Ammonius Saccas, the shadowy teacher of Plotinus, is an old and
unresolved question. And at some point ``when Zephyrinus was bishop of
Rome''---between about 198 and 217---``Adamantius, for this also was a name
of Origen,'' traveled to the capital, ``desiring, as he himself somewhere
says, to see the most ancient church of Rome,'' and returned to his work
(\work{Church History} 6.14.10--11). It is the first appearance in the record
of the nickname---``man of adamant''---by which admirers would remember him.
